% p.4 — الباب الأول في الصرف الصغير
\mspage{4}
\bab{الباب الأول في الصرف الصغير}
\begin{sarf}{|r|c|c|c|c|c|c|}
\thead{اصطلاحات} & \thead{} & \thead{وزن} & \thead{مثال} & \thead{مثال} & \thead{مثال} & \thead{ضمائر}\\\hline
ماضي معروف & 1 & فَعَلَ & فَتَحَ & جَعَلَ & دَفَعَ & هُوَ\\
مضارع معروف & 2 & يَفْعَلُ & يَفْتَحُ & يَجْعَلُ & يَدْفَعُ & هُمَا\\
مصدر معروف & 3 & فَعْلًا & فَتْحًا & جَعْلًا & دَفْعًا & هُمْ\\
اسم فاعل & 4 & فَاعِلٌ & فَاتِحٌ & جَاعِلٌ & دَافِعٌ & هِيَ\\
ماضي مجهول & 5 & فُعِلَ & فُتِحَ & جُعِلَ & دُفِعَ & هُمَا\\
مضارع مجهول & 6 & يُفْعَلُ & يُفْتَحُ & يُجْعَلُ & يُدْفَعُ & هُنَّ\\
مصدر مجهول & 7 & فَعْلًا & فَتْحًا & جَعْلًا & دَفْعًا & أَنْتَ\\
اسم مفعول & 8 & مَفْعُوْلٌ & مَفْتُوْحٌ & مَجْعُوْلٌ & مَدْفُوْعٌ & أَنْتُمَا\\
أمر حاضر & 9 & اِفْعَلْ & اِفْتَحْ & اِجْعَلْ & اِدْفَعْ & أَنْتُمْ\\
نهي حاضر & 10 & لَا تَفْعَلْ & لَا تَفْتَحْ & لَا تَجْعَلْ & لَا تَدْفَعْ & أَنْتِ\\
\twoline{اسم ظرف}{\footnotesize زمان ومكان} & 11 & مَفْعَلٌ & مَفْتَحٌ & مَجْعَلٌ & مَدْفَعٌ & أَنْتُمَا\\
اسم آلة & 12 & مِفْعَلٌ & مِفْتَحٌ & مِجْعَلٌ & مِدْفَعٌ & أَنْتُنَّ\\
اسم تفضيل & 13 & أَفْعَلُ & أَفْتَحُ & أَجْعَلُ & أَدْفَعُ & أَنَا\\
اسم مرة & 14 & فَعْلَةٌ & فَتْحَةٌ & جَعْلَةٌ & دَفْعَةٌ & نَحْنُ\\
اسم مبالغة & 15 & فَعَّالٌ & فَتَّاحٌ & جَعَّالٌ & دَفَّاعٌ & ---\\\hline
\multicolumn{2}{c}{\thead{اصطلاح}} & \thead{اصطلاح} & \thead{تمرين} & \thead{تمرين} & \thead{تمرين} & \thead{تمرين}\\\hline
\multicolumn{2}{c}{مُذَكَّرٌ} & جَمْعٌ & رَفَعَ & طَحَنَ & خَلَعَ & قَرَأَ\\
\multicolumn{2}{c}{مُؤَنَّثٌ} & غَائِبٌ & مَدَحَ & قَطَعَ & قَلَعَ & قَرَعَ\\
\multicolumn{2}{c}{وَاحِدٌ} & مُخَاطَبٌ & زَرَعَ & طَبَخَ & سَبَحَ & طَمَعَ\\
\multicolumn{2}{c}{تَثْنِيَةٌ} & مُتَكَلِّمٌ & جَمَعَ & & مَنَعَ & قَنَعَ\\
\end{sarf}
% Layout note: the serial numbers 1–15 sit as a narrow unlabelled column between اصطلاحات and وزن;
% in the lower (تمرين) block that column does not exist, so the اصطلاح cells span both.
